% !TEX root = ../main.tex
\chapter{Pendahuluan}

\section{Latar Belakang}
Pemeliharaan pembangkit listrik merupakan aktivitas penting untuk mempertahankan keandalan alat, kontinuitas operasi, dan ketersediaan sistem produksi energi kedepannya. Sektor ketenagalistrikan saat ini mengalami transformasi digital yang mendorong pemanfaatan data operasional untuk mendukung pengambilan keputusan teknis. Solusi berbasis AI dalam sektor tenaga listrik digunakan pada berbagai area seperti optimisasi sistem, deteksi gangguan, dan predictive maintenance \cite{franki2023ai_power_sector}. Selain itu, penerapan AI pada bidang operation and maintenance juga menjadi salah satu fokus penting karena berkaitan langsung dengan pemeliharaan aset pembangkitan dan infrastruktur tenaga listrik \cite{franki2023ai_power_sector}. Dengan demikian, data pemeliharaan pembangkit listrik memiliki peran strategis sebagai sumber informasi untuk memahami kondisi aktual pekerjaan perawatan dan mendukung evaluasi proses operasional.

\noindent Data pemeliharaan pada sistem ketenagalistrikan tidak hanya berisi catatan administratif, tetapi juga mempresentasikan aktivitas teknis yang berkaitan dengan kondisi peralatan, riwayat pekerjaan, dan proses perawatan. Data produksi dapat mencakup area \textit{electricity}, \textit{maintenance}, dan \textit{security}, serta berkaitan dengan sistem seperti \textit{enterprise asset management}, \textit{performance analyst}, \textit{real-time database}, \textit{safety supervision}, dan \textit{equipment reliability system} \cite{ge2017electric_power_data_mining}. Selain itu, pekerjaan \textit{overhaul} atau pemeliharaan dijelaskan sebagai bagian penting dalam operasi sistem tenaga listrik karena melibatkan status peralatan, analisis kerusakan, serta pengendalian proses pekerjaan melalui \textit{order} perbaikan \cite{ge2017electric_power_data_mining}.
Oleh karena itu, catatan pemeliharaan pembangkit listrik relevan untuk diolah sebagai \textit{event log} agar urutan aktivitas aktual dapat dianalisis menggunakan pendekatan \textit{process mining}.

Untuk mengatasi tantangan transparansi proses tersebut, disiplin ilmu \textit{Process Mining}, khususnya teknik \textit{process discovery}, menawarkan pendekatan berbasis data untuk merekonstruksi alur kerja aktual secara otomatis dari \textit{event logs} sistem informasi \cite{vanderAalst2016_process_mining}. Teknik ini bertujuan memvisualisasikan bagaimana proses benar-benar berjalan, bukan sekadar bagaimana proses tersebut didesain. Namun, penerapan algoritma \textit{discovery} konvensional pada lingkungan pemeliharaan pembangkit listrik yang kompleks sering kali terkendala oleh tingginya \textit{noise} dan variasi pencatatan aktivitas teknis. Akibatnya, model proses yang dihasilkan cenderung rumit (\textit{spaghetti models}) dan sulit dipahami, karena algoritma standar bekerja secara ``buta'' terhadap konteks aktivitas pemeliharaan yang tercermin dalam catatan operasional dan deskripsi pekerjaan \cite{augusto2019automated}.

Perkembangan \textit{Generative Artificial Intelligence} (GenAI), khususnya \textit{Large Language Models} (LLM), membuka cara baru untuk mempertemukan data log dengan pengetahuan domain. Sejumlah studi memanfaatkan LLM sebagai perantara yang mengekstrak aturan dari deskripsi proses tekstual untuk memandu algoritma \textit{process discovery} agar menghasilkan model yang lebih presisi \cite{norouzifar2024}. Arah lain menelaah \textit{event abstraction}, yaitu penyeragaman label aktivitas dari tingkat rinci ke tingkat yang lebih bermakna, dan memperlihatkan bahwa langkah ini memperbaiki keterbacaan model proses \cite{tax2016event}. Kajian terbaru menegaskan model bahasa mampu menangani tugas \textit{process mining} yang bertumpu pada pemahaman makna aktivitas, termasuk \textit{process discovery} itu sendiri \cite{rebmann2025semantics}. Peluang inilah yang diadopsi pada penelitian ini, yakni menstandarkan label aktivitas pemeliharaan sebelum tahap \textit{process discovery} dijalankan, alih-alih sekadar menyaring \textit{raw data}.

Berdasarkan urgensi untuk meningkatkan efisiensi operasional pada proses pemeliharaan pembangkit listrik, penelitian ini berfokus pada pemanfaatan GenAI untuk mendukung \textit{process discovery}. Integrasi kecerdasan buatan generatif dengan \textit{process mining} diharapkan mampu menghasilkan model alur kerja pemeliharaan yang tidak hanya akurat secara data, tetapi juga lebih sederhana, presisi, dan mudah diinterpretasikan melalui penyeragaman label aktivitas, sehingga dapat mendukung evaluasi proses serta pengambilan keputusan teknik berbasis data.

\section{Perumusan Masalah}

Penelitian ini difokuskan pada pemanfaatan data log yang tersedia. Permasalahan yang akan dijawab adalah

\begin{enumerate}
    \item Bagaimana Generative AI dapat dimanfaatkan untuk menstandarkan ragam label aktivitas pemeliharaan menjadi sejumlah kategori yang konsisten?
    \item Bagaimana bentuk model proses pemeliharaan yang dihasilkan dari perbandingan \textit{process discovery} pada \textit{event log as-is} dan \textit{event log enriched}?
    \item Apakah penyeragaman label aktivitas berbasis LLM menghasilkan model proses yang lebih ringkas dibandingkan metode konvensional tanpa anotasi?
\end{enumerate}


\section{Batasan Masalah}

Untuk menjaga fokus dan konsistensi analisis, penelitian ini dibatasi pada beberapa aspek sebagai berikut. Objek penelitian hanya mencakup satu unit Pembangkit Listrik Tenaga Diesel (PLTD) Titikuning. Sumber data dibatasi pada dokumen \textit{job card} pekerjaan pemeliharaan dan struktur \textit{Work Breakdown Structure} (WBS) dari kegiatan \textit{overhaul} unit PLTD Titikuning. Dengan demikian, temuan yang dihasilkan tidak dimaksudkan untuk digeneralisasikan ke unit pembangkit lain maupun ke jenis pemeliharaan yang berbeda.

Pada tahap anotasi, penelitian ini memanfaatkan model \texttt{gemini-3.1-flash-lite} dengan \textit{temperature} 0,1 agar keluaran bersifat konsisten. \textit{Event} dengan \textit{confidence} di bawah 0,45 diklasifikasikan anomali agar dapat ditinjau ulang. Pada tahap \textit{process discovery}, algoritma yang digunakan dibatasi pada \textit{Inductive Miner infrequent} (IMf) dengan \textit{noise threshold} 0,2 dan diimplementasikan melalui PM4Py. Penstandaran label aktivitas dibatasi pada 12 kategori yang ditetapkan penulis, sehingga penyeragaman dilakukan dalam kerangka taksonomi yang telah ditentukan, bukan secara bebas tanpa batas.


\section{Rencana Kegiatan}

Penelitian ini dilaksanakan melalui tahapan-tahapan sebagai berikut:

\begin{enumerate}
  \item \textbf{Studi Literatur dan Pemahaman Alur}
  \begin{itemize}
    \item Mempelajari konsep dasar \textit{Process Mining}, khususnya cara algoritma bekerja dalam memetakan data log menjadi model alur proses.
    \item Mendalami kemampuan \textit{Generative AI}, khususnya \textit{Large Language Model} (LLM), dalam memahami teks bahasa manusia dan mengubahnya menjadi informasi yang terstruktur.
    \item Mengkaji karakteristik alur kerja pemeliharaan pembangkit listrik, khususnya pekerjaan \textit{overhaul} pada unit Pembangkit Listrik Tenaga Diesel (PLTD), untuk memahami istilah teknis yang muncul dalam data (misalnya pembersihan, pemeriksaan, pengukuran, dan penggantian komponen), sehingga penulis memiliki bekal pemahaman awal terhadap data.
  \end{itemize}

  \item \textbf{Pengumpulan dan Pra-pemrosesan Data (Data Preprocessing)}
  \begin{itemize}
    \item \textbf{Pengumpulan Data:} Mengumpulkan sampel data catatan pemeliharaan (\textit{maintenance log}) dari sistem informasi yang tersedia. Fokus utama adalah memperoleh data yang memuat setidaknya tiga unsur, yaitu nomor pekerjaan, deskripsi pekerjaan (teks), dan waktu pengerjaan.
    \item \textbf{Penyesuaian Format Data:} Mempelajari struktur data yang diperoleh, lalu melakukan pembersihan data (\textit{data cleaning}) untuk memilah data yang layak olah. Tahap ini meliputi penghapusan data kosong dan penyesuaian format tanggal agar dapat dibaca oleh perangkat lunak.
    \item \textbf{Pembentukan Event Log:} Mengubah \textit{raw data} menjadi format yang dapat dibaca oleh algoritma \textit{Process Mining}. Aktivitas majemuk pada satu baris pekerjaan dipecah menjadi beberapa \textit{event} melalui \textit{compound activity splitting}, dan setiap pekerjaan direpresentasikan dengan penanda waktu mulai serta selesai agar durasi pekerjaan dapat terlihat.
  \end{itemize}

  \item \textbf{Perancangan Metode dan Prompt Engineering}
  \begin{itemize}
    \item Merancang serangkaian perintah (\textit{prompt}) untuk \textit{Large Language Model} (LLM) agar dapat membaca kolom deskripsi pekerjaan yang berupa teks bebas dan memetakannya ke dalam kategori tindakan berbasis taksonomi.
    \item Menguji \textit{prompt} pada sebagian data untuk memastikan model mampu memahami konteks pekerjaan pemeliharaan pembangkit listrik dengan tepat.
  \end{itemize}

  \item \textbf{Implementasi Sistem dan Eksperimen}
  \begin{itemize}
    \item \textbf{Pengembangan Kode:} Menulis program menggunakan Python untuk menghubungkan data dengan algoritma \textit{Process Mining}.
    \item \textbf{Pelaksanaan Eksperimen:} Menjalankan dua skenario pemodelan:
    \begin{enumerate}
      \item \textit{Skenario as-is:} membangun model proses langsung dari \textit{raw data} tanpa anotasi AI.
      \item \textit{Skenario enriched:} membangun model proses dengan memanfaatkan hasil anotasi kategori tindakan oleh LLM.
    \end{enumerate}
  \end{itemize}

  \item \textbf{Evaluasi Model}
  \begin{itemize}
    \item \textbf{Analisis Perbandingan:} Membandingkan model proses dari kedua skenario untuk menilai mana yang lebih ringkas, tidak ruwet (mengurangi efek \textit{spaghetti model}), dan lebih mudah dipahami.
    \item \textbf{Pengukuran Kualitas:} Mengukur kualitas model menggunakan metrik \textit{conformance}, yaitu \textit{fitness}, \textit{precision}, \textit{generalization}, dan \textit{simplicity}, untuk memberikan penilaian objektif terhadap model yang dihasilkan.
  \end{itemize}

  \item \textbf{Analisis Hasil dan Penyusunan Laporan}
  \begin{itemize}
    \item Menganalisis apakah pemanfaatan AI benar-benar membantu menyederhanakan pemetaan proses pemeliharaan pembangkit listrik.
    \item Mendokumentasikan seluruh proses, tantangan selama pengolahan data, dan hasil temuan ke dalam Laporan Tugas Akhir.
  \end{itemize}
\end{enumerate}

\clearpage
\section{Jadwal Kegiatan}

Penelitian ini diselesaikan dalam enam bulan dengan rincian kegiatan seperti pada Tabel~\ref{tab:jadwalKegiatan}.

\begin{table}[H]
  \centering
  \caption{Jadwal Kegiatan Penelitian Tugas Akhir}
  \label{tab:jadwalKegiatan}
  \footnotesize
  \renewcommand{\arraystretch}{1.5}
  {\setlength{\tabcolsep}{3pt}
  \newcommand{\x}{\cellcolor{blue!30}}
  \begin{tabularx}{\textwidth}{|c|X|*{6}{p{0.25cm}p{0.25cm}p{0.25cm}p{0.25cm}|}}
  \hline
  \multirow{2}{*}{\textbf{No}} &
  \multirow{2}{*}{\textbf{Kegiatan}} &
  \multicolumn{24}{c|}{\textbf{Bulan ke-}} \\
  \cline{3-26}
   & &
  \multicolumn{4}{c|}{\textbf{1}} &
  \multicolumn{4}{c|}{\textbf{2}} &
  \multicolumn{4}{c|}{\textbf{3}} &
  \multicolumn{4}{c|}{\textbf{4}} &
  \multicolumn{4}{c|}{\textbf{5}} &
  \multicolumn{4}{c|}{\textbf{6}} \\
  \hline
  1 & Studi Literatur            & \x & \x & \x &    &    &    &    &    &    &    &    &    &    &    &    &    &    &    &    &    &    &    &    &    \\ \hline
  2 & Pengumpulan Data           & \x & \x & \x & \x & \x & \x &    &    &    &    &    &    &    &    &    &    &    &    &    &    &    &    &    &    \\ \hline
  3 & Pra-pemrosesan Data        &    &    &    &    & \x & \x & \x & \x & \x & \x &    &    &    &    &    &    &    &    &    &    &    &    &    &    \\ \hline
  4 & Perancangan Mekanisme Berbasis Generative AI/LLM &    &    &    &    &    &    &    &    & \x & \x & \x & \x & \x & \x &    &    &    &    &    &    &    &    &    &    \\ \hline
  5 & Integrasi dengan Algoritma Process Discovery     &    &    &    &    &    &    &    &    &    &    &    &    & \x & \x & \x & \x & \x & \x &    &    &    &    &    &    \\ \hline
  6 & Evaluasi Model Proses      &    &    &    &    &    &    &    &    &    &    &    &    &    &    &    &    & \x & \x & \x & \x & \x & \x & \x &    \\ \hline
  7 & Penyusunan Laporan         &    &    &    &    &    &    &    &    &    &    &    &    &    &    &    &    &    &    &    &    &    & \x & \x & \x \\ \hline
  \end{tabularx}
  }
\end{table}

\newpage

